\documentclass[11pt,a4paper]{article}

\usepackage[utf8]{inputenc}
\usepackage[sr-cyrl]{babel}
\usepackage[T2A]{fontenc}

\usepackage{graphicx}
\usepackage{amsmath, amsthm, amssymb}
\usepackage{xcolor}
\usepackage{wrapfig}
\usepackage{geometry}
\geometry{a4paper, margin=2.5cm}


\usepackage{hyperref}
\hypersetup{
    colorlinks=true,
    linkcolor=black,     
    citecolor=blue,      
    urlcolor=blue        
}

\newtheorem{definicija}{Дефиниција}[section]
\newtheorem{teorema}{Теорема}[section]
\newtheorem{lema}{Лема}[section]

\title{\textbf{\Large МАТЕМАТИЧКИ ФАКУЛТЕТ}\\ \vspace{0.5cm} \Huge 5G/6G мреже и рачунарска оптимизација}
\author{\textbf{Студент:} Тамара Спасеновски (106/2025) \\ \textbf{Професор:} Данијела Симић}
\date{Београд, 01. 09. 2026.}

\begin{document}

\maketitle
\thispagestyle{empty}
\clearpage

\pagenumbering{roman}
\tableofcontents
\clearpage
\pagenumbering{arabic}

\section{Увод}
Од првих генерација мобилне телефоније, чија је основна намена била пренос гласа, до данашњих напредних мрежа, мобилне технологије су значајно утицале на савремено друштво. Савемене $5G$ и долазеће $6G$ мреже захтевају изузетно сложен математички апарат и алгоритме рачунаРСке оптимизације како би ефикасно управљале мрежним ресурсима.

\section{Развој и карактеристике мобилних мрежа}

\subsection{Еволуција мобилних генерација}
\begin{itemize}
    \item \textbf{1G i 2G:} Аналогни пренос гласа, затим дигитализација и увођење $SMS$ порука.
    \item \textbf{3G i 4G:} Појава мобилног интернета, а затим $IP$ протокола са брзинамa до $100\text{ Mbps}$.
    \item \textbf{5G i 6G:} Прелазак на масовну комуникацију машине-ка-машини и екстремне перформансе.
\end{itemize}

\subsection{Поређење перформанси}
У табели испод дато је паралелно поређење свих кључних карактеристика.

\begin{table}[h!]
\centering
\caption{Поређење $4G$, $5G$ и $6G$ перформанси}
\vspace{0.2cm}
\begin{tabular}{|l|c|c|c|}
\hline
\textbf{Карактеристика} & \textbf{4G} & \textbf{5G} & \textbf{6G} \\ \hline
\textbf{Максимална брзина} & $1\text{ Gbps}$ & $20\text{ Gbps}$ & \textcolor{purple}{\textbf{1 Tbps}} \\ \hline
\textbf{Кашњење (латенција)} & $50\text{ ms}$ & $< 1\text{ ms}$ & \textcolor{blue}{\textbf{$< 0.1\text{ ms}$}} \\ \hline
\textbf{Густина уређаја ($/\text{km}^2$)} & $2.000$ & $1.000.000$ & $> 10.000.000$ \\ \hline
\textbf{Поузданост} & $99.9\%$ & $99.999\%$ & $99.99999\%$ \\ \hline
\end{tabular}
\end{table}

\section{Оптимизациони проблеми и математички апарат}

\subsection{Математичка формулација проблеми расподеле ресурса}

\begin{wrapfigure}{l}{0.4\textwidth}
  \centering
  \includegraphics[width=0.38\textwidth]{iot.jpg}
  \caption{Колосалан број $IoT$ уређаја у $5G/6G$ мрежама.}
\end{wrapfigure}

За успешно управљање великим бројем $IoT$ уређаја, као што је илустровано на слици, морамо дефинисати математичке моделе оптимизације.

\begin{definicija}[Проблем расподеле фреквенција]
Проблем оптималне расподеле $M$ радио-канала на $N$ корисника уз максимизацију укупног протока дефинише се као:
\begin{equation}
\max \sum_{i=1}^{N} \sum_{j=1}^{M} R_{ij} x_{ij}
\end{equation}
где је $x_{ij} \in \{0, 1\}$ индикатор алокације, а $R_{ij}$ остварена брзина корисника $i$ на каналу $j$.
\end{definicija}

\begin{teorema}[Сжетост проблема алокације]
Проблем оптималне алокације ресурса је $NP$-тежак за општи случај са великим бројем корисника $N \gg 10^3$.
\end{teorema}

\begin{lema}[Минимизација потрошње енергије]
Укупна потрошња $K$ базних станица минимизује се функцијом:
\begin{equation}
\min \sum_{k=1}^{K} P_k \quad \text{уз услов квалитета сигнала } \text{SINR}_k \ge \gamma_{\min}
\end{equation}
\end{lema}

\section{Методе оптимизације и машинско учење}
За решавање наведених проблема примењују се следећи приступи:
\begin{enumerate}
    \item \textbf{\textcolor{darkgray}{Конвексна оптимизација:}} Користи се када функционали имају гарантовани глобални екстремум.
    \item \textbf{\textcolor{blue}{Учење поткрепљивањем (Reinforcement Learning):}}
    \begin{enumerate}
        \item Агент бира акцију у мрежи.
        \item Окружење враћа нумеричку награду за квалитет одлуке.
        \item Агент ажурира своју интерну стратегију у реалном времену.
    \end{enumerate}
\end{enumerate}

\section{Закључак}
\textit{\textbf{\textcolor{purple}{5G и 6G мреже}}} представљају темељ будућег дигиталног друштва \cite{ref2}. Примена савремених рачунарских и математичких метода омогућава да се комплексни мрежни ресурси организују аутономно и ефикасно.

\begin{thebibliography}{9}
\bibitem{ref1} Dahlman, E., Parkvall, S., \& Skold, J. (2020). \textit{5G NR: The Next Generation Wireless Access Technology}. Academic Press.
\bibitem{ref2} Zhang, Z., et al. (2019). \textit{6G Wireless Networks: Vision, Requirements, Architecture, and Key Technologies}. IEEE Vehicular Technology Magazine.
\end{thebibliography}

\end{document}